\chapter{Introducción}

El análisis exploratorio de datos (Exploratory Data Analysis, EDA) constituye una etapa fundamental dentro del proceso de análisis de datos, ya que permite comprender la estructura del conjunto de datos, evaluar su calidad e identificar patrones, tendencias y posibles anomalías antes de aplicar técnicas estadísticas o modelos analíticos más avanzados.

En el presente caso de estudio se analiza un conjunto de datos correspondiente al sector bancario mediante herramientas de Python orientadas al análisis descriptivo y la visualización de información. Para ello se emplean métricas de calidad de datos, estadísticas descriptivas y representaciones gráficas que facilitan la interpretación de los datos.

El propósito de este informe es documentar de forma estructurada el proceso de análisis realizado, presentando tanto los resultados obtenidos como las principales conclusiones y recomendaciones derivadas del estudio.